% !PRINTED 377
\weberpage{377}
Conversely, if the congruence (3) is satisfied, one may set
$D = n^2 - 4\,ml$ and obtains a form of discriminant $D$,
$(m, n, l)$, by which $m$ is properly representable.

1. \gesperrt{The necessary and sufficient condition for a number $m$ coprime to $Q$ to be properly representable by a form of discriminant $D$ therefore consists in this, that the congruence}

\weq{4}{x^2 \equiv D \pmod{4m}}

\gesperrt{be solvable.}

From a form $(a,b,c)$ one obtains a family of \gesperrt{equivalent
forms}:

\weq{5}{(a,\ b+2\lambda a,\ C),}

in which $\lambda$ is an arbitrary integer and $C$ is determined by
the equation

\[ D = (b+2\lambda a)^2 - 4aC \]

and turns out to be equal to $a\lambda^2+b\lambda+c$. The system (5)
is called a \gesperrt{family of parallel forms}.

If the congruence (4) is satisfied, then also

\[ (x+2\lambda m)^2 \equiv D \pmod{4m}, \]

and we shall regard the whole family of roots of (4) which are
congruent with respect to the modulus $2m$ as \gesperrt{one root}. In
this sense let the number of roots of (4), which is accordingly
certainly finite, be denoted by $\psi(D,m)$.

2. \gesperrt{Every root of (4) gives rise to a family of parallel forms by which the number $m$ is properly representable.}

The number $\psi(D,m)$ is easily determined. If first
$m'$, $m''$ are relatively prime, then

\weq{6}{\psi(D,m'),\ \psi(D,m'') = \psi(D,m'm'').} % [sic] German source prints a comma rather than a multiplication sign.

For if

\weq{7}{x'^2 \equiv D \pmod{4m'}, \quad x''^2 \equiv D \pmod{4m''},}

and

\weq{8}{\begin{aligned} x &\equiv x' \pmod{2m'}, \\ &\equiv x'' \pmod{2m''}, \end{aligned}}

then $(x^2-D):4$ is divisible by $m'$ and by $m''$, hence also by
$m'm''$, and the congruence

\weq{9}{x^2 \equiv D \pmod{4m'm''}}

is satisfied. Conversely, every root of the congruence (9) gives one
root of each of the two congruences (7). Thus one obtains all
%CONT>>
